\chapter{Stores and global state}
\label{ch:stores}

\begin{aside}
Some state belongs to a widget. Other state is global: the
logged-in user, the current theme, the open documents. \maya{}
calls these stores, and gives them lifetimes that match an
application, not a widget tree.
\end{aside}

\section{What a store is}

A store is a bundle of signals exposed via a typed handle.
Multiple widgets can read and write the same store.

\begin{lstlisting}
struct ThemeStore {
    Signal<std::string> name;
    Signal<Color> accent;
};

auto theme = ctx.make_store<ThemeStore>(
    "dark", "#FF6432"
);
\end{lstlisting}

The store lives at root scope --- accessible from any widget.

\section{Dependency injection}

A widget accepts the store as a constructor argument:

\begin{lstlisting}
auto title_bar = [&] {
    auto name = ctx.get(theme->name);
    return text("Theme: " + name);
};
\end{lstlisting}

This is dependency injection by passing references. No
global singletons.

\section{Context propagation}

For deep trees, passing stores everywhere is verbose. The
context system propagates:

\begin{lstlisting}
ctx.provide(theme);          // at root
auto t = ctx.consume<ThemeStore>();  // deep child
\end{lstlisting}

Equivalent to React's Context API.

\section{Actions and reducers}

For complex stores, encapsulate writes behind actions:

\begin{lstlisting}
class DocumentStore {
    Signal<std::vector<Doc>> docs;
public:
    void add(Doc d) {
        auto v = ctx.get(docs);
        v.push_back(d);
        ctx.set(docs, v);
    }
    void remove(int i) { /* ... */ }
};
\end{lstlisting}

Widgets call actions; the store maintains invariants.

\section{Selectors}

A selector is a memoised derivation:

\begin{lstlisting}
auto unread_count = memo([&] {
    auto docs = ctx.get(docs);
    return std::count_if(docs.begin(), docs.end(),
                         [](auto& d) { return !d.read; });
});
\end{lstlisting}

Widgets that need just the count subscribe to the selector,
not the full docs list. Re-renders only when count changes.

\section{Persistence}

Stores often need to persist across runs. \maya{} provides:

\begin{lstlisting}
ctx.persist(theme, "theme.json");
\end{lstlisting}

The runtime serialises the store on change and reloads on
startup.

\section{Time travel}

A store with history allows undo/redo:

\begin{lstlisting}
auto history = HistoricStore(store);
history.undo();
\end{lstlisting}

Each action pushes a snapshot; undo pops.

\section{Recap}

\begin{recap}
\begin{itemize}[leftmargin=*]
  \item Stores: bundles of signals at root scope.
  \item Inject via constructor or context.
  \item Actions encapsulate writes.
  \item Selectors memoise derivations.
  \item Persist to file; history for undo.
\end{itemize}
\end{recap}

\section*{Workshop}

\begin{enumerate}[leftmargin=*]
  \item Build a theme store and switch themes at runtime.
  \item Add a document store with actions for add/remove.
  \item Persist your store to disk and reload on startup.
\end{enumerate}
